\documentclass{article}

\usepackage{../../apxproof}
\usepackage{cleveref}

\newtheoremrep{theorem}{Theorem}
\newtheoremrep*{lemma}{Lemma}

\begin{document}

\section{Main text}

\begin{theoremrep}
  \label[theorem]{th:opt}%
  A theorem labeled with an explicit counter type. It has an equation:
  \begin{equation}
    \label[equation]{eq:opt}
    1=1
  \end{equation}
  and one labeled without an explicit type:
  \begin{equation}
    \label{eq:noopt}
    2=2
  \end{equation}
\end{theoremrep}

\begin{proof}
  Deferred proof, referring to \cref{eq:opt} and \cref{eq:noopt}.
\end{proof}

\begin{lemmarep}
  \label[lemma]{lem:opt}%
  A starred repeated theorem, also labeled with a counter type.
\end{lemmarep}

\begin{proof}
  Deferred proof of \cref{lem:opt}.
\end{proof}

\begin{toappendix}
  \label[section]{app:opt}%
  Appendix material, labeled with a counter type.
\end{toappendix}

References to the main text: \cref{th:opt}, \cref{eq:opt}, \cref{eq:noopt},
\cref{lem:opt}, \cref{app:opt}.

References to the appendix copies of the equations: \cref{eq:opt-apx},
\cref{eq:noopt-apx}.

\end{document}
